% !TeX root = document.tex
\subsubsection{Базова реалізація: прямий (нативний) метод}

Практична реалізація розробленого алгоритму починається з програмного представлення цифрового зображення як двовимірної дискретної матриці пікселів $F$ розмірністю $W \times H$. Граничні просторові індекси цієї матриці позначаються як $n$ та $m$, причому фізична ширина зображення становить $W = n + 1$, а висота — $H = m + 1$. Кожен елемент матриці з просторовими індексами $(i, j)$ зберігає значення кольору $f_{i,j}$ у форматі RGB. Задачею масштабування є побудова нової матриці зображення із цільовими розмірами $W_2 \times H_2$, де нова ширина $W_2 = N + 1$, а нова висота $H_2 = M + 1$ ($N$ та $M$ відповідають максимальним індексам цільової просторової сітки).

Першим етапом розробки стала програмна імплементація прямого (нативного) алгоритму двовимірної барицентричної інтерполяції, що представлений у коді функцією \texttt{processDirectBarycentricZoom}. Процес геометричного перетворення починається з відображення просторових кординат пікселів обох матриць у неперервний двовимірний інтервал $[-1, 1] \times [-1, 1]$. Для вихідного масиву даних просторова інтерполяційна сітка формується за законом розподілу екстремумів поліномів Чебишева другого роду:
\begin{equation}\label{eq:chebyshev_nodes}
	x_i = - \cos\left(\frac{\pi i}{n}\right), \quad y_j = - \cos\left(\frac{\pi j}{m}\right), \quad \text{де } i = \overline{0, n}, \quad j = \overline{0, m}.
\end{equation}

Так само просторові координати пікселів цільового зображення $(X, Y)$ переводяться у математичний діапазон $[-1, 1]$ за тим же правилом. У результаті ми отримуємо нормалізовані координати $(x, y)$, для яких безпосередньо обчислюється нове значення кольору:
\begin{equation}\label{eq:target_coordinates}
	x = - \cos\left(\frac{\pi X}{N}\right), \quad y = - \cos\left(\frac{\pi Y}{M}\right), \quad \text{де } X = \overline{0, N} \quad Y = \overline{0, M}.
\end{equation}

Для обчислення значення функції кольору $C(x, y)$ кожного результуючого пікселя використовується двовимірна друга барицентрична формула, яка враховує внесок усіх без винятку пікселів вхідного зображення:
\begin{equation}\label{eq:barycentric_2d_formula}
	C(x, y) = \frac{\sum\limits_{i=0}^{n} \vphantom{\sum}^{\prime} \sum\limits_{j=0}^{m} \vphantom{\sum}^{\prime} \frac{(-1)^{i+j}}{(x - x_i)(y - y_j)} f_{i,j}}{\sum\limits_{i=0}^{n} \vphantom{\sum}^{\prime} \sum\limits_{j=0}^{m} \vphantom{\sum}^{\prime} \frac{(-1)^{i+j}}{(x - x_i)(y - y_j)}}.
\end{equation}

З точки зору програмної архітектури, функція \texttt{processDirectBarycentricZoom} приймає об'єкт класу \texttt{QImage}, виконує попіксельне вилучення колірних компонент за допомогою низькорівневого доступу до пам'яті, ініціалізує результуючий об'єкт \texttt{QImage} заданих цільових розмірів і запускає обчислювальний процес. Важливою умовою стабільності програми є обробка ситуацій, коли поточні координати точки $(x, y)$ збігаються з вузлами вихідної просторової сітки Чебишева. У таких випадках знаменники у виразі \eqref{eq:barycentric_2d_formula} перетворюються на нуль, викликаючи критичні обчислювальні помилки. Для запобігання цьому у програмному коді реалізовано умовне розгалуження, логіка якого повністю спирається на математичне обґрунтування граничних випадків, що наведено вище. Якщо фіксується точний або частковий просторовий збіг, алгоритм автоматично блокує розрахунок загальної двовимірної форми та зводить операцію до одновимірного аналогу або безпосереднього копіювання значення кольору $f_{i,j}$.

Незважаючи на високу математичну точність, проведений аналіз обчислювальної ефективності прямого методу виявив його повну практичну непридатність для систем реального часу. Архітектура нативного алгоритму передбачає наявність чотирьох глибоко вкладених циклів: зовнішні цикли ітеруються по просторових координатах цільової матриці ($Y$ від $0$ до $M$, $X$ від $0$ до $N$), а всередині них для кожного кроку запускаються два вкладені цикли повного перебору вихідних точок ($j$ від $0$ до $m$, $i$ від $0$ до $n$). 

Така структура зумовлює вкрай високу асимптотичну складність алгоритму, яка оцінюється як $\mathcal{O}(W \cdot H \cdot W_2 \cdot H_2)$. Наприклад, якщо здійснюється просторове збільшення відносно невеликого графічного масиву розміром $500 \times 500$ пікселів у два рази (до розміру $1000 \times 1000$), процесору необхідно виконати:
\begin{equation}\label{eq:complexity_calc}
	1000 \times 1000 \times 500 \times 500 = 2.5 \cdot 10^{11}
\end{equation}
тобто 250 мільярдів ітерацій, кожна з яких містить операції ділення та множення. Таке колосальне навантаження викликає велике завантаження одного обчислювального потоку процесора, призводить до значних затримок (порядку кількох хвилин на один кадр) та спричиняє ефект тимчасового блокування (зависання) головного вікна графічного інтерфейсу користувача програми.

Реалізований прямий алгоритм на сітці Чебишева виконав свою основну роль у проекті — він підтвердив абсолютну коректность розробленого математичного апарату та виступив у ролі програмного еталона. Проте виявлені критичні обмеження швидкодії та нераціональне використання ресурсів процесора довели необхідність перебудови архітектури обчислень, що привело до розробки оптимізованої моделі.